\section{Iterative Deployment Improves Planning Skills in LLMs}
\label{sec:mechanism}

We study \textbf{Iterative Deployment}, a training mechanism that leads Large Language Models (LLMs) to bootstrap their planning capabilities without requiring external expert demonstrations or additional teacher models. The core intuition is that a model can improve by learning from the simple tasks it has already successfully solved, effectively using its own valid outputs as training data for subsequent generations (provided a reliable curation/verification mechanism is available).

We hypothesize that LLMs improve their planning abilities by simply fine-tuning
over their own curated traces. In other words, they can bootstrap by first
solving easy tasks themselves, and then being fine-tuned using the traces of
these solved tasks. Later generations can then start solving larger tasks, and
these new traces can be used so future generations solve even larger
ones. Repeating this process many times can then gradually improve their planning skills. For
example, if the current generation of a model can solve Sokoban problems with
one single box, then it can learn from its own traces to solve problems with two
boxes, and then with three boxes, etc. So by exploiting its own current
capabilities, the model can improve and solve harder tasks in future
generations.

Formally, let $M_n$ denote the model at generation $n$, parameterized by $\theta_n$. We assume access to a dataset of planning tasks \emph{without} solutions $\mathcal{D}_{tasks}$ and a deterministic external validator $V(x, y)$ which returns true if trace $y$ is a valid solution for task $x$, and false otherwise, with an ability to measure solution efficiency. This validator can be seen as a correction mechanism for reasoning tasks, or even simply as a proxy for user preferences. We evaluate model performance on  $\mathcal{D}_{tasks}$ as a test-set as well. This is because we are interested to find how iterative deployment improves the model's ability to solve longer tasks with access only to its own previously curated solutions. The iterative deployment process then proceeds as follows:

\begin{itemize}
    \item \textbf{Deployment and Trace Collection:} In each iteration $n$, we prompt the current model $M_n$ to solve the tasks in $\mathcal{D}_{tasks}$. For each task input $x$, the model generates a trace $y$ according to its policy $\pi_{\theta_n}(y|x)$. This trace includes the chain-of-thought and the tentative solution generated by the model. This emulates a standard deployment scenario where the model interacts with users or environments.
    \item \textbf{Validation:} The generated traces are passed to the external validator $V$. We filter the outputs to retain only the subset of valid traces, $\mathcal{D}_{valid}^{(n)} = \{(x, y) |V(x, y) = \text{True}\}$. Invalid plans, which certainly constitute the majority of outputs in early generations, are discarded.
    \item \textbf{Curation and Aggregation:} To prevent catastrophic forgetting, reduce model collapse, and further improve generalization, we aggregate the valid traces from the current generation with those from all previous generations. The training dataset for the next step is $\mathcal{T}_{n+1} = \bigcup_{i=0}^{n} \mathcal{D}_{valid}^{(i)}$.
    \item During this aggregation, we apply a second curation step: a \textbf{selection mechanism} to ensure data quality. If multiple valid traces exist for the same task (e.g., from different generations), we retain only the highest-quality solution. In our experiments, quality is defined by plan efficiency -- we select the trace with the shortest plan length, breaking ties by selecting the one with fewer reasoning tokens, but in principle other task-specific metrics can be used.
    \item \textbf{Supervised Fine-Tuning (SFT):} Finally, we produce the next generation $M_{n+1}$ by fine-tuning $M_n$ on the curated dataset $\mathcal{T}_{n+1}$ using the standard supervised learning objective (next token prediction).
\end{itemize}

The curation and aggregation phase is important here, as the likelihood of a
user making an interaction with an LLM publicly available is not uniform, but
depends on the nature of the interaction and the users' intention. For instance,
when a user interacts with an LLM to solve a coding task, they are more likely
to integrate a response that solves their tasks into their codebase. They are
even more likely to use it if the solution is particularly elegant or simple,
akin to our selection mechanism. This creates an effective curation mechanism
controlled by the user's revealed preference (which is not necessarily their
stated preference). This is a key difference to the assumptions used in the study of model collapse \cite{shumailov-et-al-nature2024}.
As LLMs are trained on more and more recursively generated data, their
performance can deteriorate and the models eventually collapse. However, curation might delay or prevent
this.
In this work, we study the effect of this additional assumption compared to the model collapse assumptions, and
argue for the importance of understanding its impact on future model
generations.

Figure~\ref{fig:pipeline} illustrates how one iteration of this process
works. When the $n$-th generation of an LLM is deployed, it is prompted by their users. The produced traces are
filtered by an external \emph{validator}, and those judged \emph{valid traces}
are used to fine-tune the $n+1$-th generation of the model.
Note that we are \emph{not} building a learning curriculum ourselves, but rather the LLM together with the validator are building one: we prompt the model with all tasks from the test set $\mathcal{D}_{tasks}$, discard the invalid traces, add the valid ones to the train set, and then fine-tune the next generation.
